% Proposed B.94.  The retained-deck extension is separate from B.93's E-span.
\statement{B.94}{Proved}{All retained-fiber deck translates form a free rank-64 module with the same four-source Tate image}
\label{stmt:B94}
Retain B.91--B.93. Let $T=F(y_1,\ldots,y_5)$ be the actual retained
generic-fiber group in the corrected paired marking, with involution
\begin{equation}\label{eq:B94-beta}
 \beta(y_1)=y_1^{-1},\quad
 \beta(y_2)=y_4,\quad\beta(y_3)=y_5,\quad
 \beta(y_4)=y_2,\quad\beta(y_5)=y_3.
\end{equation}
The common free peripheral basis of $R$ is indexed by
$x_{g,i}$, where $g\in T$ is a retained Schreier word and
$i\in\{2,3,5,6,7,8,10,11\}$ is a deleted paired slot.
Put $\alpha_t(x_{g,i})=x_{tg,i}$ and define
\begin{equation}\label{eq:B94-G}
 \begin{aligned}
 G&=(R\times R)\rtimes T,\qquad \Omega=\mathbb Z[G],\\
 (u,v,t)(x,y,s)&=
       \bigl(u\alpha_t(x),v\alpha_{\beta(t)}(y),ts\bigr),\\
 F(u,v,t)&=(v,u,\beta(t)).
 \end{aligned}
\end{equation}
This is the actual generic-fiber quotient acting on $P$ after
abelianizing the mixed commutator kernel. In particular it includes
all retained-fiber deck translates, without a length bound.
For the $64$ actual columns of B.93, set
\[
                 L_{\mathrm{ret}}=\sum_{j=0}^{63}\Omega c_j\subset P.
\]
Then $L_{\mathrm{ret}}$ is $F$-stable and
\begin{equation}\label{eq:B94-result}
 \begin{gathered}
 L_{\mathrm{ret}}\cong\Omega^{64},\qquad
 H^1(C_2,L_{\mathrm{ret}})\cong\mathbb F_2[R]^4,
 \qquad H^2(C_2,L_{\mathrm{ret}})=0,\\
 \operatorname{im}\bigl(H^1(C_2,L_{\mathrm{ret}})
                  \longrightarrow H^1(C_2,P)\bigr)
       =J_4=\sum_{i=1}^4\mathbb F_2[R](\rho_i-1)\subseteq J_I.
 \end{gathered}
\end{equation}
The four $\rho_i$ are the same actual first-sheet source elements
as in B.93. Thus every integral anti-chain in this entire
retained-deck span has zero $\kappa$ value modulo the actual
source ideal, interpreted by B.93's intrinsic refinement.

The old product-pencil base operators and their arbitrary iterates
are not included in this coefficient enlargement. Their additional
orbit sum in the full transported presentation remains untreated.
No global $\kappa$ vanishing, projected source equality, nonlinear
fixed lifting, compact comparison, or invertibility of the selected
matrix is asserted.

\begin{proof}
\emph{B.94.1. Actual normal forms and a free coefficient-row module.}
In the actual generic fiber, kill the retained filling kernel and
obtain the free group $T$. The peripheral kernel before imposing
mixed node relations is $R_A*R_B$, with the actual Schreier bases
of B.91. For an $A$ peripheral, a retained whisker $g$ gives the
common label $(g,i)$. For a $B$ peripheral, a retained whisker $h$
gives the common label $(\beta(h),i)$. This makes $F$ exchange
the two factors with identical common labels. Conjugation by
$t\in T$ therefore acts by $\alpha_t$ on the first factor and
by $\alpha_{\beta(t)}$ on the second. After killing their mixed
commutators, the actual quotient has the unique normal form
\eqref{eq:B94-G}. Multiplication, inverses and the involution
follow from these actions and $\beta^2=1$.

A product-tree cell of $P$ has two vertices and two peripheral
labels, which we write
\[
                         (u,v;(g,i),(h,j)).
\]
Its unique coefficient-row normal form is
\begin{equation}\label{eq:B94-row-normal}
 (u,v;(g,i),(h,j))=(u,v,g)\,\epsilon_{i,j,z},
                  \qquad z=\beta(g)^{-1}h,
\end{equation}
where $\epsilon_{i,j,z}$ has vertices $(1,1)$ and labels
$(1,i),(z,j)$. Indeed, applying the retained coefficient $g$
changes the two labels to $(g,i),(\beta(g)z,j)$, and applying
$(u,v)$ supplies the vertices. Conversely, the first label
determines $g$, after which the second determines $z$ and the
vertices determine $u,v$. This proves that $P$ is a free left
$\Omega$-module on the rows $(i,j,z)$; it is not merely a
surjection from a proposed free row module.

Write $d_t=(1,1,t)$. Literal negative interchange of the two
product-tree edges gives
\begin{equation}\label{eq:B94-row-F}
 F(\epsilon_{i,j,z})
       =-d_z\epsilon_{j,i,\beta(z)^{-1}}.
\end{equation}
Together with the coefficient action in \eqref{eq:B94-G}, this
recovers the original $F$ action on every cell. Applying it twice
returns the original row and coefficient. The exact conversion
\path{work_CE/code/convert_symmetric_retained_deck_module.py}
checks \eqref{eq:B94-row-normal} by inverse expansion on all
$2680$ terms of the $64$ actual columns, and checks the complete
row involution. There are $282$ distinct used rows. The formulas
prove the normal form for every row and every retained word;
the finite check verifies the actual input conversion.

\emph{B.94.2. The actual $64$ monomial involution identities.}
Let $\nu$ be a physical node and $\mu$ its physical $F$-mate.
The actual matching calculation retains the full based conjugator,
its residual words and its retained $T$ head. With the $A$ branch
ordered first, it gives the complete column equation
\begin{equation}\label{eq:B94-mate}
 F(c_\nu)=-\delta_\nu a_\nu c_\mu,
 \qquad a_\nu\in G,\quad\delta_\nu\in\{1,-1\},
 \qquad F(a_\nu)a_\mu=1,
 \quad\delta_\nu\delta_\mu=1.
\end{equation}
All $64$ equations are checked as full $G$-linear chains in
\path{work_CE/next_symmetric_nodes/symmetric_retained_deck_module.json}.
No nontrivial retained head is discarded or replaced by an
$E$ coefficient. The equations make $L_{\mathrm{ret}}$ $F$-stable.
They define a semilinear involution on the free domain on its
$64$ columns. Proving this domain maps injectively to $P$ is
a separate step, supplied next.

\emph{B.94.3. A separated integral augmentation filtration.}
Let $I_E$ be the integral augmentation ideal of $\mathbb Z[E]$
and put $\mathfrak a=I_E\Omega$. Retained conjugation preserves
$I_E$, so normal form in the semidirect product gives
\begin{equation}\label{eq:B94-powers}
 \mathfrak a^n=I_E^n\Omega
       =\bigoplus_{t\in T}I_E^n d_t.
\end{equation}
We prove directly that these powers have zero intersection. The
power-series construction is classical
\citep[\S2, Satz I, pp.~262--263]{Magnus1935}; the characteristic-two
and finite-quotient separation arguments needed here are proved below.

First, a free group of any rank is residually a finite $2$-group.
Only finitely many generators occur in a chosen nontrivial word,
so kill the others. Write its reduced maximal-block expression as
$x_{i_1}^{a_1}\cdots x_{i_r}^{a_r}$, where $a_j\ne0$ and
$i_j\ne i_{j+1}$. In the noncommutative formal power series
algebra over $\mathbb F_2$, send $x_i$ to $1+X_i$.
If $a=2^v b$ with $b$ odd, the first nonzero positive-degree
term of $(1+X)^a-1$ has degree $2^v$ and coefficient $1$.
This holds also for negative $a$, by formal inversion, or by
$(1+X)^a=(1+X^{2^v})^b$ in characteristic two.
Put $d_j=2^{v_2(a_j)}$. The monomial
\[
                     X_{i_1}^{d_1}\cdots X_{i_r}^{d_r}
\]
has exactly $r$ successive runs of variable names. Omitting
the positive-degree term from any block gives fewer than $r$
runs; choosing from every block gives this monomial only with
the displayed degrees. Its coefficient is therefore $1$.
The image of the word is nontrivial already in a finite
degree truncation. The group $1+\mathfrak m$ of that truncated
algebra is a finite $2$-group: its filtration by
$1+\mathfrak m^j$ has elementary abelian $2$-group quotients.
This supplies a finite $2$-group quotient detecting the word.
Products of finitely many such quotients separate any finite
set of distinct words. Applying this to the two free factors
shows that $E=R\times R$ has the same finite-support separation
property.

For a finite $2$-group $V$, the augmentation ideal of
$\mathbb F_2[V]$ is nilpotent. Here is the elementary induction.
Choose a central element $z$ of order two. The central ideal
generated by $z-1$ has square zero, and quotienting by it gives
$\mathbb F_2[V/\langle z\rangle]$. The quotient augmentation
ideal is nilpotent by induction on $|V|$, so the original is
nilpotent as well. If its $N$th power is zero modulo two,
the augmentation ideal in $(\mathbb Z/2^k)[V]$ has $N$th
power contained in $2(\mathbb Z/2^k)[V]$, and hence its
$Nk$th power is zero.

Now let $a\ne0$ have finite support in $\mathbb Z[E]$.
Choose a finite $2$-group quotient of $E$ separating its
support, and choose $k$ so that one nonzero integer coefficient
survives modulo $2^k$. The image of $a$ in
$(\mathbb Z/2^k)[V]$ is nonzero. Nilpotence of the target
augmentation ideal proves that $a$ cannot lie in every
$I_E^n$. Thus
\begin{equation}\label{eq:B94-separated}
              \bigcap_{n\geq1}I_E^n=0,
              \qquad\bigcap_{n\geq1}\mathfrak a^n=0.
\end{equation}
The second assertion follows coefficient by coefficient from
\eqref{eq:B94-powers}. No finite quotient invariant under
all of $T$ is needed. The coefficient reduction modulo $2^k$
is essential; nilpotence over the integral group ring of a
finite group was not asserted.

\emph{B.94.4. Injectivity of the actual column map.}
Apply the residual augmentation
$\Omega\to\mathbb Z[T]$ to the normalized matrix. Each
of the $64$ columns has exactly one nonzero augmented entry,
a signed retained word, and the $64$ surviving rows are distinct.
Their ordered physical slot pairs are exactly the $8$ by $8$
product of the slot set in the statement. These facts are
verified in
\path{work_CE/next_symmetric_nodes/symmetric_G_augmentation_pivots.json};
the producer separately reconstructs the selected $64$ by $64$
matrix and verifies its diagonal form.

Choose those rows in the corresponding column order. The actual
matrix is $D+M$, where $D$ has diagonal entries
$u_i=\pm d_{t_i}\in\Omega^\times$ and every entry of $M$
belongs to $\mathfrak a$. Suppose
$\sum_i a_i c_i=0$, with arbitrary $a_i\in\Omega$ acting
on the left. In the $i$th selected row the equation is
\[
                       a_i u_i+\sum_j a_j m_{ji}=0.
\]
Modulo $\mathfrak a$, the diagonal units imply $a_i\in
\mathfrak a$ for every $i$. If all $a_j\in\mathfrak a^n$,
then the same ordered equation gives
\[
                 a_i=-\sum_j a_jm_{ji}u_i^{-1}
                          \in\mathfrak a^{n+1}.
\]
Induction and \eqref{eq:B94-separated} force every $a_i=0$.
The $64$ actual columns are therefore $\Omega$-independent,
and by definition they generate $L_{\mathrm{ret}}$. This
proves the free rank in \eqref{eq:B94-result}. The argument
proves injectivity only. It neither inverts $D+M$ nor treats
an element with unit augmentation as a unit of $\Omega$.

\emph{B.94.5. Fixed coefficient points and the two retained cone classes.}
There are $28$ exchanged pairs of physical nodes and eight
fixed physical nodes. In each fixed case $\delta_\nu=1$,
so \eqref{eq:B94-mate} has the form $F(c)=-a c$, where
$a=(u,v,t)$ and $F(a)a=1$. Its underlying coefficient-basis
involution is
\begin{equation}\label{eq:B94-coefficient-orbit}
                              x\longmapsto F(x)a.
\end{equation}
It has a fixed point if and only if some $w\in T$ satisfies
\begin{equation}\label{eq:B94-retained-half}
                              \beta(w)t=w.
\end{equation}
Necessity follows by projection to $T$. For sufficiency, the
cocycle identity gives
$\beta(t)t=1$ and $v\alpha_{\beta(t)}(u)=1$.
If \eqref{eq:B94-retained-half} holds, direct multiplication
in \eqref{eq:B94-G} shows that
\[
                              x=(\alpha_{\beta(w)}(u),1,w)
                              \quad\text{satisfies}\quad F(x)a=x.
\]
Indeed $\beta(w)=w\beta(t)$, so its second-coordinate equation
is $\alpha_{\beta(w)}(u)\alpha_w(v)=1$ and its other two
coordinates agree literally.

There is an explicit free-product presentation
\begin{equation}\label{eq:B94-retained-orbifold}
 T\rtimes\langle\tau\rangle
    =\langle\tau\mid\tau^2=1\rangle*
      \langle\tau_1\mid\tau_1^2=1\rangle*F(y_2,y_3),
 \qquad y_1=\tau_1\tau,
\end{equation}
where $y_4=\tau y_2\tau$, $y_5=\tau y_3\tau$ and conjugation
by $\tau$ on $T$ is $\beta$. Substitution in both directions
verifies the presentation. Equation \eqref{eq:B94-retained-half}
is equivalent to $\tau t$ being conjugate to $\tau$ by an
orientation-even element. The two order-two factors have
distinct conjugacy classes by reduced free-product normal form.
Also the centralizer of $\tau$ in this free product is
$\langle\tau\rangle$. Its intersection with $T$ is trivial,
so
\begin{equation}\label{eq:B94-fixed-G}
                    T^\beta=1,
                    \qquad G^F=\{(r,r,1):r\in R\}.
\end{equation}
If $x_0$ is one fixed point of \eqref{eq:B94-coefficient-orbit},
all its fixed points are exactly the left torsor $G^F x_0$:
the quotient $x x_0^{-1}$ is fixed by $F$, and the converse
follows by multiplication. This fixes the side of the resulting
source ideal.

For completeness, the eight actual retained tests are as follows.
The table gives $t$ and the single factor remaining after cyclic
reduction of $\tau t$ in \eqref{eq:B94-retained-orbifold}.
The full reducing prefixes and residual coefficients are retained
in the exact converted data.
\begin{center}
\begin{tabular}{ccc}
Node&Retained word $t$&Cone core\\\hline
$(0,0)$&$y_1$&$\tau_1$\\
$(0,1)$&$1$&$\tau$\\
$(8,0)$&$y_2y_3y_1^{-1}y_5^{-1}y_4^{-1}$&$\tau_1$\\
$(8,1)$&$1$&$\tau$\\
$(12,0)$&$y_2y_1^{-1}y_4^{-1}$&$\tau_1$\\
$(12,1)$&$1$&$\tau$\\
$(16,0)$&$y_1$&$\tau_1$\\
$(16,1)$&$1$&$\tau$
\end{tabular}
\end{center}
The four other-cone cases have no fixed coefficient point,
regardless of their residual coordinates. For the four neutral
cases, the actual coefficient is $(b,b^{-1},1)$, with
$b=\rho_1,\rho_2,1,1$ respectively. We choose the explicit
fixed coefficient point
\[
                              x_0=(1,b^{-1},1),
\]
whose anti-chain is precisely the corresponding $v_i$ of
B.93.3. Thus the previous actual
source normalization is retained exactly.

\emph{B.94.6. Cohomology and the actual source image.}
Independence in B.94.4 makes the monomial relations
\eqref{eq:B94-mate} the actual action on a free image basis.
Each exchanged pair is induced as a $C_2$-module. In a fixed
physical summand with no fixed coefficient point, every
coefficient-basis orbit has size two, so it too is a direct
sum of induced signed pairs and has no positive-degree
cohomology. This is why fixing a physical node alone does
not produce a Tate class.

For a neutral fixed summand, its fixed coefficient points
are the left diagonal $R$-torsor in \eqref{eq:B94-fixed-G}.
Each carries the negative sign, while all other points lie
in induced pairs. The integral sign-module calculation from
B.93 therefore gives one copy of $\mathbb F_2[R]$ in
$H^1$ and zero in $H^2$ for each of the four neutral summands.
All chains have finite support, so the orbitwise calculation
applies to their direct sum without an infinite-support
choice or division by two.

The four distinguished anti-chains are exactly B.93's $v_i$.
Their Tate values are the actual $\rho_i-1$, and the fixed
coefficient torsor gives their diagonal left translates
$s(\rho_i-1)$. Thus the full cohomology image is exactly
$J_4$, proving all of \eqref{eq:B94-result}. In particular,
adjoining all retained-deck translates contributes no further
class modulo $J_I$.

\emph{B.94.7. Exact evidence and the remaining orbit extension.}
\begingroup\raggedright
The actual normal-form conversion, all $64$ coefficient
identities, their ordered double-$F$ checks and all eight
cone tests are in
\path{work_CE/next_symmetric_nodes/symmetric_retained_deck_module.json},
with producer
\path{work_CE/code/convert_symmetric_retained_deck_module.py}
and genuine adjacent output. The separate $64$ diagonal
augmentation units are in
\path{work_CE/next_symmetric_nodes/symmetric_G_augmentation_pivots.json},
with producer
\path{work_CE/code/check_symmetric_G_augmentation_pivots.py}
and genuine adjacent output. Both certificates bind their
actual inputs by SHA-256 and retain the complete words and
coefficients. The separated-filtration proof above is the
general injectivity argument; the finite computation verifies
its exact actual-matrix hypothesis.
The independent replay
\path{work_DE/next_marginal_source/check_G_span_independent.py}
reconstructs all $64$ normalized columns, the $2680$ terms,
all mate and cycle relations, the selected diagonal minor
and all eight cone classes from the actual inputs; its
genuine adjacent output records each successful check.\par
\endgroup

The coefficient group $G$ is the whole generic fiber modulo
the mixed commutator kernel. That kernel acts trivially on
its abelianization, so $\Omega$ coefficients already include
all generic-fiber conjugations of these columns. The further
node transports in B.91 come from the old product-pencil
base paths and their action on this fiber data. Their orbit
sum is not replaced by $L_{\mathrm{ret}}$ here, nor is
$L_{\mathrm{ret}}$ asserted to be invariant under those
operators. B.93's induced refinement places every computed
anti-chain in the exact $\kappa$ framework, but the full
old-base sum and the subsequent nonlinear and compact gates
remain open.
\end{proof}
